\documentclass[10pt,a4paper]{article}
\usepackage{amsmath,amssymb,amsthm,graphicx}
\usepackage{amsmath,amssymb,amsthm,hyperref}
\newtheorem{nttheorem}{Theorem}
\newtheorem{ntdefinition}[nttheorem]{Definition}
\newtheorem{ntproposition}[nttheorem]{Proposition}
\newtheorem{ntlemma}[nttheorem]{Lemma}
\newtheorem*{ntclaim}{Claim}
\newcommand{\DTIME}{\operatorname{DTIME}}
\newcommand{\SIZE}{\operatorname{SIZE}}
\newcommand{\poly}{\operatorname{poly}}
\title{Circuit lower bounds from parametric-honest or natural Turing NP-hardness of MCSP}
\author{Michael Saks, Rahul Santhanam}
\date{doi:10.4230/LIPIcs.CCC.2020.26; CC BY3.0}
\begin{document}
\maketitle
\noindent\textit{Editorial note.} Michael Saks and Rahul Santhanam. CCC2020 article26,DOI10.4230/LIPIcs.CCC.2020.26. CC BY3.0.\par\medskip
\noindent\textit{Editorial note.} Faithful editable transcription of the official CCC2020 CC BY3.0 publication, article26 pp26:1--26:13. One structured-Turing-hardness-to-circuit-lower-bound theme: introductory Theorems1/2, complete Section2 definitions/prerequisites3--6, complete Section3 Theorems7/8 author proofs, and complete Section4 Lemma9/Theorem10 full author proofs, including both simulations, padding construction, interval Claim and its full argument, robustly-often dichotomy, small hard-instance extraction and hard-query concatenation. All24 original bibliography entries retained. Parametric honest and natural are separate incomparable author conditions, both required only on correctly answered prefixes; their exact definitions and separate proof routes retained. No claim of a consequence for unrestricted/general Turing reductions, whose question remains open in original Section5. External standard diagonalization/Karp-Lipton/robust-simulation prerequisites and Gutfreund-Shaltiel-TaShma result remain at original author citation boundary without recursive proof. Original printed exponents/indices and wording retained without mathematical repair, including interval induction nprime\_(j+1)=r(n\_j), Claim strict/non-strict signs, first-disjunct T(n)=poly(r(n))n\textasciicircum{}s1(r(n)), and parameter s2 in failure-formula construction. Authors retain mathematical responsibility. No native source was provided by official page; formal PDF/text faithfully transcribed. No figures required. Root handles compilation/database/publication.\par\medskip
\noindent\textit{Allowed standard prerequisites (editorial).}\begin{itemize}
\item Robust simulation and time diagonalization. Original Proposition3 citing[11,Theorem12] and Proposition4 citing[13] retained with exact time conditions.
\item Karp-Lipton theorem for E and hard-instance finding. Original author Karp-Lipton/Meyer[19] and Gutfreund-Shaltiel-TaShma[12] uses and stated three-instance guarantee retained; no recursion into prior proofs.
\end{itemize}\medskip
\setcounter{nttheorem}{0}
\begin{nttheorem}[Author Theorem1]
If MCSP is NP-hard under parametric honest Turing reductions, then $\mathrm E\not\subseteq\SIZE(\poly)$.
\end{nttheorem}
\begin{nttheorem}[Author Theorem2]
If MCSP is NP-hard under natural Turing reductions, then $\mathrm E\not\subseteq\SIZE(\poly)$.
\end{nttheorem}

\section*{2 Preliminaries}
\subsection*{2.1 Basic Notions}
We refer to the book by Arora and Barak~\cite{ref9} for definitions of standard complexity classes.
Given a circuit class $\mathcal C$, we use $\mathcal C[s(n)]$ to refer to the class of functions computed by $\mathcal C$-circuits of size $s(n)$.
Given a language $L\subseteq\{0,1\}^*$, $\operatorname{co}\text{-}L$ denotes the complement of $L$. Given language $L$ and $n\in\mathbb N$, $L_n$ is used to refer to $L\cap\{0,1\}^n$.
Given a set $S\subseteq\mathbb N$ and languages $L$ and $L'$, we say $L$ agrees with $L'$ on $S$ if for all $n\in S$, $L_n=L'_n$.
A time-constructible function $T:\mathbb N\to\mathbb N$ is a function such that there is a Turing machine transducer $M$, which for each $n$, on input $1^n$ halts with output $T(n)$ within $O(T(n))$ steps.
We need a generalization of the notion of robustly-often simulation from~\cite{ref11}. Given $g:\mathbb N\to\mathbb N$, a set $S\subseteq\mathbb N$ is called $g$-robust if for each constant $k$, there is $m\in\mathbb N$ such that for all $n$ such that $m\leq n\leq g(m)^k$, $n\in S$. Note that any $g$-robust set is also $\poly(g)$-robust.
Given a function $g$, a language $L$ and a complexity class $\mathcal C$, we say $L$ is $g$-robustly often in $\mathcal C$ if there is $L'\in\mathcal C$ for which there is a $g$-robust set $S$ such that $L$ agrees with $L'$ on $S$. Given function $g$ and complexity classes $\mathcal C$ and $\mathcal C'$, we say that $\mathcal C$ is $g$-robustly often in $\mathcal C'$ if for all languages $L\in\mathcal C$, $L$ is $g$-robustly often in $\mathcal C'$. We use the term ``robustly often'' as a synonym for $g$-robustly often with $g$ the identity function.
The proposition below is a parameterized version of Theorem12 in~\cite{ref11}.
\setcounter{nttheorem}{2}
\begin{ntproposition}[Author Proposition3]
Let $T$ be a time-constructible function such that $T(n)\geq n$ for all $n$. If SAT is $T(\poly(n))$-robustly often in $\DTIME(T)$, then $\Sigma_2\mathrm P$ is robustly often in $\DTIME(T(\poly(T(\poly(n)))))$.
\end{ntproposition}
The following is a simple application of standard diagonalization~\cite{ref13}.
\begin{ntproposition}[Author Proposition4]
Let $T$ be a time-construtible function such that $T=o(2^n)$. Then $\mathrm E$ is not robustly often in $\DTIME(T)$.
\end{ntproposition}
\subsection*{2.2 Resource-Bounded Kolmogorov Complexity}
We study various notions of resource-bounded Kolmogorov complexity, and associated decision problems.
Fix a universal Turing machine $U$. The $Kt$ complexity of a string is defined as follows:
\[Kt(x)=\min\{|p|+\log(t)\mid U(p)\text{ halts and outputs }x\text{ within }t\text{ steps}\}.\]
In the MKtP problem, the instance is a string $x$ together with a parameter $s$, and the question is whether $Kt(x)\leq s$.
It is known that MKtP is complete for $\mathrm E$ under polynomial-size truth-table reductions~\cite{ref2}, but completeness under uniform polynomial-time reductions is unknown.
In the MCSP problem, the instance is a string $x$ together with a parameter $s$, and the question is whether $x$ is the truth table of a Boolean function with circuit complexity at most $s$.
MCSP is easily seen to be in NP, but it is unknown whether MCSP is NP-hard. A brute-force search strategy of trying all circuits $C$ of size at most $s$ and checking if any of them computes the function with truth table $x$ can easily be implemented to run in time $\poly(|x|)s^{O(s)}$.
\subsection*{2.3 Reductions}
We will work with constrained notions of polynomial-time Turing reduction.
\begin{ntdefinition}[Author Definition5]
A polynomial-time Turing reduction from a language $L$ to MCSP or MKtP is called parametric honest if there is a constant $\epsilon>0$ such that for every $x\in\{0,1\}^*$, every query made by the reduction on $x$ has parameter bounded below by $|x|^\epsilon$.
\end{ntdefinition}
Note that parametric honesty is not implied by the standard notion of honesty for polynomial-time reductions, which only requires the output length to be bounded below by some polynomial in the input length. Implicit in the definition is that the lower bound on parameters of queries is only required to hold when previous queries are answered correctly; when queries are answered wrongly, ``all bets are off''.
\begin{ntdefinition}[Author Definition6]
A polynomial-time Turing reduction from a language $L$ to MCSP or MKtP is called natural if for any input $x\in\{0,1\}^*$, the parameter of any query made by the reduction on $x$ depends only on $|x|$.
\end{ntdefinition}
This notion of naturalness generalizes the notion used by Kabanets and Cai for $m$-reductions to Turing reductions. In fact, it is slightly weaker than their notion when restricted to $m$-reductions, since it does not constrain the length of the output, merely the parameter. As with the definition of parametric honesty, the condition on parameters is only required to hold when previous queries are answered correctly.
Naturalness is incomparable to parametric honesty--parametric honesty allows the parameter to vary but restricts its range of variation, while naturalness allows the parameter to have a large range of variation but insists that it is the same for all queries made on any input of a given length.

\section*{3 Parametric Honest Reductions}
Our first result unconditionally rules out hardness of MKtP for $\mathrm E$ under parametric honest Turing reductions. For many results on hardness of MCSP and consequences thereof, MKtP is a ``test case'' where analogous results can be shown unconditionally and where the ideas are often simpler~\cite{ref8,ref16}.
\begin{nttheorem}[Author Theorem7]
MKtP is not hard for $\mathrm E$ under parametric honest Turing reductions.
\end{nttheorem}
\begin{proof}
Let $L$ be a tally set that is in $\mathrm E$ but not in $\mathrm P$. The existence of such a tally set $L$ follows by a standard diagonalization argument.
Now suppose, for the purpose of contradiction, that MKtP is hard for $\mathrm E$ under parametric honest Turing reductions. This implies that there is an oracle Turing machine $M$ running in time at most $cn^k$, where $c$ and $k\geq1$ are constants, such that $M$ decides $L$ correctly with oracle access to MKtP, and moreover there is a constant $\epsilon>0$ such that each query made by $M$ on an input of length $n$ has parameter at least $n^\epsilon$. We show how to use $M$ to decide $L\in\mathrm P$, contrary to the assumption on $L$.
Consider the following polynomial-time machine $A$ that attempts to decide $L$. Given input $0^n$, $A$ runs the oracle machine $M$, answering each query made by $M$ with ``yes''. When the simulation of $M$ is complete, $A$ accepts iff $M$ accepts. Since $M$ is polynomial-time, so is $A$. We show that $A$ correctly decides $L$ for large enough $n$.
Given input $0^n$, $M$ makes some sequence of oracle queries $q_1,q_2\ldots q_t$ during the simulation by $A$, where $0\leq t\leq cn^k$. We show that for each $i$, $q_i$ has $Kt$ complexity $O(\log(n))$. Observe that $q_i$ can be generated by the universal machine $U$ implicit in the definition of $Kt$ complexity in time $\poly(n)$ given descriptions of $i,n$ and the oracle machine $M$: it simply needs to run $M$ on $0^n$ answering the first $i-1$ queries made by $M$ with ``yes'', and then output the $i$'th query. The description length required is $O(\log(n))$ since $i\leq cn^k$, and moreover the time taken by $U$ is polynomial in $n$, hence using the definition of $Kt$ complexity, we have an $O(\log(n))$ upper bound on the $Kt$ complexity of $q_i$. Let $C$ be a constant such that the $Kt$ complexity of each $q_i$ is bounded above by $C\log(n)$.
By the assumption that the reduction is parametric honest, the parameter of $q_i$ is at least $n^\epsilon$, which for large enough $n$ is greater than $C\log(n)$. Hence for large enough $n$ and for each $i$, it follows that $A$'s assumed answer for the oracle query $q_i$ is correct for each $i$, and hence that $A$ outputs the same answer as $M$ would given the true oracle MKtP. Since $M$ is an oracle Turing machine correctly implementing a reduction from $L$ to MKtP, $A$ is a polynomial-time Turing machine correctly deciding $L$ for large enough $n$, contradicting the assumption on $L$.
\end{proof}
We now re-state and prove Theorem1.
\begin{nttheorem}[Author Theorem8]
If MCSP is NP-hard under parametric honest Turing reductions, then $\mathrm E\not\subseteq\SIZE(\poly)$.
\end{nttheorem}
\begin{proof}
Suppose that MCSP is NP-hard under parametric honest Turing reductions. Let $M$ be a polynomial-time oracle Turing reduction implementing a parametric honest reduction from SAT to MCSP, and let $\epsilon$ be a constant such that the parameter of any query made by $M$ on any SAT instance of length $x$ is bounded below by $|x|^\epsilon$.
Consider the following polynomial-time machine $A$ which attempts to decide SAT. On an input $x$ of SAT, $A$ runs the oracle Turing machine $M$. Assume wlog that all queries asked by $M$ have length a power of two--any queries for which this is not the case may be eliminated by assuming the answer is `no'. For every query asked by $M$ with length a power of two, $A$ assumes the answer is ``yes'', and continues the simulation of $M$. When the simulation concludes, $A$ accepts iff $M$ accepts.
Clearly $A$ runs in polynomial time. If $A$ decides SAT correctly, then we have that $\mathrm{NP}=\mathrm P$. Hence the Polynomial Hierarchy collapses to $\mathrm P$. But this implies that $\mathrm E\not\subseteq\SIZE(\poly)$; if the contrary were true, then $\mathrm E$ would collapse to the Polynomial Hierarchy using the Karp--Lipton theorem for $\mathrm E$, and hence to $\mathrm P$, contradicting the hierarchy theorem for deterministic time.
Suppose now that $\mathrm{NP}\ne\mathrm P$. In particular, this means that there are infinitely many instances on which $A$ does not decide SAT correctly. If not, we could hardcode the finitely many instances on which $A$ fails to compute SAT into a new polynomial-time machine $A'$ which runs in polynomial time and solves SAT correctly on all instances, contradicting the assumption that $\mathrm{NP}\ne\mathrm P$.
Now, we can use the main result of Gutfreund, Shaltiel and Ta-Shma. They show that if $\mathrm{NP}\ne\mathrm P$, then for any polynomial-time algorithm $A$ that fails to compute SAT, there is a polynomial-time Turing machine $B$ and a constant $d$, such on input $1^n$, $B$ outputs three instances $x_1,x_2,x_3$ each of length $n$ or length $n^d$ such that for infinitely many $n$, $A$ fails to decide SAT correctly on at least one of these three instances.
We show how to use $B$ to argue again that $\mathrm E\not\subseteq\SIZE(\poly)$. Consider the three instances $x_1,x_2,x_3$ output by $B$ on input $1^n$. For each $i$, $1\leq i\leq3$, let $Q_i$ be the concatenation of all queries made by $M$ when $M$ is simulated by $A$ on input $x_i$. Let $Y$ be the concatenation of $Q_1,Q_2,Q_3$. Note that $|Y|=\poly(n)$, since $M$ is a polynomial-time oracle Turing machine. Also note that each query made by $M$ on an input of length $n$ has parameter at least $n^\epsilon$, by the parametric honesty of the reduction, as long as all previous queries were answered correctly. Each such query was assumed by $A$ to have a ``yes'' answer. It could not be the case that all such answers were correct for all queries made by $M$ on inputs $x_1,x_2,x_3$--if they were, then $A$ would correctly solve SAT on $x_1,x_2,x_3$, using the fact that $M$ implements a correct polynomial-time Turing reduction from SAT.
Thus we get that there is a first query made by $M$ on one of $x_1,x_2,x_3$ that has parameter at least $n^\epsilon$ and is wrongly answered ``yes''. Fix such a query $y$. Since $M$ implements a Turing reduction from SAT to MCSP, $y$ is the truth table of a Boolean function on $O(\log(n))$ bits that requires circuits of size $n^\epsilon$. We have that $y$ is a substring of $Y$, hence $Y$ must also require circuits of size $\Omega(n^\epsilon)$. We are using here the simple fact that if $y$ is a substring (with length a power of two) of the truth table $Y$ of a Boolean function with circuits of size $s$, then $y$ is the truth table of a Boolean function with circuits of size $O(s)$.
Hence, for infinitely many $n$, on input $1^n$, we can compute the truth table $Y$ of a Boolean function on $O(\log(n))$ bits such that the function requires circuits of size $n^\epsilon$, i.e., of exponental size in the length of the input to the Boolean function. This implies $\mathrm E\not\subseteq\SIZE(\poly)$.
\end{proof}

\section*{4 Natural Reductions}
We require a technical lemma allowing us to use indirect diagonalization based on robustly-often simulations.
\begin{ntlemma}[Author Lemma9]
Let $T:\mathbb N\to\mathbb N$ be a time-constructible function such that $T(\poly(T(\poly(n))))=o(2^n)$. If SAT is $T(\poly(n))$-robustly often in time $T(n)$, then $\mathrm E\not\subseteq\SIZE(\poly)$.
\end{ntlemma}
\begin{proof}
Suppose SAT is $T(\poly(n))$-robustly often in time $T(n)$. By Proposition3, we have that $\Sigma_2\mathrm P$ is robustly often in time $T(\poly(T(\poly(n))))$. Assume for the sake of contradiction that $\mathrm E\subseteq\SIZE(\poly)$. The Karp--Lipton theorem for $\mathrm E$, credited to Meyer~\cite{ref19}, implies that $\mathrm E$ collapses to $\Sigma_2\mathrm P$. Note that if $T$ is time-constructible, $T'=T(\poly(T(\poly(n))))$ is also time-constructible. Thus we have that $\mathrm E$ is robustly often in time $T'$ for a time-constructible $T'=o(2^n)$, contradicting Proposition4.
\end{proof}
We now re-state and prove Theorem2. In the proof below, we use the fact that SAT is paddable--there is a polynomial-time computable padding function $f$ such that for each $x\in\{0,1\}^*$, and for each $m\geq|x|$, $|f(x,1^m)|=m$ and $f(x,1^m)$ is in SAT iff $x$ is in SAT. Of course, this depends on the encoding of SAT, but it suffices for us that there is some encoding for which paddability is true in the form above, and all other standard properties of SAT continue to hold.
\begin{nttheorem}[Author Theorem10]
If MCSP is NP-hard under natural Turing-reductions, then $\mathrm E\not\subseteq\SIZE(\poly)$.
\end{nttheorem}
\begin{proof}
Suppose that MCSP is NP-hard under natural Turing reductions. Let $M$ be an oracle Turing machine implementing a natural reduction from SAT to MCSP. By the definition of natural reduction, this means that there is a function $s:\mathbb N\to\mathbb N$ such that all queries of $M$ on input of size $n$ for SAT have parameter $s(n)$. If there are no queries that $M$ makes on an input of size $n$, we set $s(n)=0$.
Let $s_1(n)=\log(n)^{\log\log(n)}$, and $s_2(n)=\log(n)^{\log\log\log(n)}$. Moreover let $r(n)=n^{\log(n)}$. We define two Turing machines $M_1$ and $M_2$ that attempt to decide SAT using the oracle machine $M$.
Machine $M_1$ operates as follows. On input $x$ of size $n$, $M_1$ computes strings $x_1\ldots x_{r(n)-n}$, where $x_i=f(x,1^{n+i})$. Here $f$ is the padding function for SAT. Let $x_0=x$. $M_1$ iteratively does the following starting with $i=0$. It simulates $M$ on $x_i$. If an oracle query is made with parameter $k$, it checks if $k$ is at most $s_1(n+i)$. Note that since $M$ implements a natural reduction, $k=s(n+i)$. If $k$ is at most $s_1(n+i)$, it solves the corresponding MCSP instance by brute force search in time at most $\poly(n+i)2^{s_1(n+i)^2}$ and continues the simulation. Note that any succeeding queries must also have parameter at most $s_1(n+i)$, therefore the simulation continues to completion. $M_1$ accepts iff $M$ accepts. If no oracle queries are made, then too the simulation continues to completion, and $M_1$ returns the same answer as $M$.
In case the parameter $k$ for a query is greater than $s_1(n+i)$, $M_1$ increments $i$, and repeats the process, as long as $i\leq r(n)-n$. If $i>r(n)-n$, $M_1$ simply runs the simulation of $M$ on $x$ to completion, now answering all queries by brute-force search regardless of the parameter size.
\begin{ntclaim}
Either SAT is decidable in time $\poly(r(n))2^{s_1(r(n))^2}$, or there is an infinite sequence of disjoint intervals $I_j=[n_j,n'_j)$ of input lengths such that $n'_j=r(n_j)$ for each $j$, and moreover for each $j$, if $n\in I_j$, then $s(n)>s_1(n)$.
\end{ntclaim}
We establish the Claim. Suppose that SAT is not decidable in time $\poly(r(n))2^{s_1(r(n))^2}$. We argue that for any $n_0\in\mathbb N$, there exists an input length $n>n_0$ such that $s(n)\geq s_1(n)$, and moreover $s(m)\geq s_1(m)$ for every $n<m<r(n)$. Indeed, if this were not the case, for every $n>n_0$ one of the iterations of $M_1$ would succeed for some $i<r(n)-n$, and this would mean that the entire simulation would complete in time at most $\poly(r(n))2^{s_1(r(n))^2}$. Note that when the simulation completes, it does solve SAT correctly, as $M$ is an oracle Turing reduction correctly reducing SAT to MCSP, and any completing simulation uses correct answers to oracle queries.
Now we pick the intervals $I_j$ inductively as follows. Let $n_1$ be the least integer such that $s(m)>s_1(m)$ for every $n_1\leq m\leq r(n_1)$, and let $n'_1=r(n_1)$. Inductively, suppose we have picked $n_1\ldots n_j$. We pick $n_{j+1}$ to be the least integer such that $n_{j+1}>r(n_j)$, and moreover $s(m)>s_1(m)$ for each $n_{j+1}\leq m\leq r(n_{j+1})$. Such an integer exists by the argument of the previous paragraph. Let $n'_{j+1}=r(n_j)$, and set $I_j=[n_j,n'_j)$. This sequence of intervals is clearly disjoint, and satisfies the property in the Claim.
The first disjunct of the Claim implies that $\mathrm E\not\subseteq\SIZE(\poly)$, using Lemma9 and the fact that $T(n)=\poly(r(n))n^{s_1(r(n))}$ is time-constructible and satisfies $T(\poly(T(\poly(n))))=o(2^n)$, for the given choice of $r$ and $s_1$. Hence, in the rest of our proof, we assume that the second disjunct holds, i.e., that there is an infinite sequence of long intervals within which each input length generates queries with a large parameter.
Now we define machine $M_2$. On input $x$ of length $n$, $M_2$ uses $M$ to try to solve the search version of SAT, i.e., to find a satisfying assignment for $x$ if one exists. $M_2$ uses the standard search-to-decision reduction for SAT based on self-reducibility, simulating $M$ to answer any decision questions as described below. In order to find a satisfying assignment for $x$, this search-to-decision procedure might call the decision procedure for SAT on inputs $x'$ such that $|x'|<|x|$; in such cases, $M_2$ pads up $x'$ to length $|x|$ by using the padding function $f(x',1^{|x|})$ and runs the simulation of $M$ on $f(x',1^{|x|})$. This guarantees that all inputs on which $M$ is run during the simulation have the same length. At the end of the search-to-decision procedure, if a satisfying assignment is successfully found, $M_2$ accepts, else it rejects.
Next we describe how $M_2$ uses $M$ to decide if an input is in SAT. $M_2$ simulates the oracle machine $M$ on the input. If $M$ asks a query with parameter greater than $s_2(n)$, $M_2$ assumes the answer is ``yes'' and continues the simulation. If $M$ asks a query with parameter at most $s_2(n)$, $M_2$ solves the query by brute force search and continues the simulation. At the end of the simulation, $M_2$ accepts iff $M$ accepts.
$M_2$ always halts within time $\poly(n)2^{s_2(n)^2}$. Unlike machine $M_1$, machine $M_2$ is not guaranteed to solve SAT correctly. However, since $M_2$ has been modified to solve the search version of SAT, the only way $M_2$ can make a mistake on $x$ is to answer ``no'' when the true answer is ``yes''. We will argue that whether or not $M_2$ solves SAT correctly, circuit lower bounds follow.
We consider two cases. The first is that there is a $2^{s_2(n)^3}$-robust set $S$ on which $M_2$ solves SAT correctly. In this case, we can apply Lemma9 with $T(n)=\poly(n)2^{s_2(n)^2}$, noting that $T$ is time-constructible and that $T(\poly(n))<2^{s_2(n)^3}$ for large enough $n$. Thus we get that $\mathrm E\not\subseteq\SIZE(\poly)$ in this case.
The remaining case is that there is no $2^{s_2(n)^3}$-robust set $S$ on which $M_2$ solves SAT correctly. In this case, we apply the main idea of~\cite{ref12}, using in addition our assumption that there is an infinite sequence $\{I_j\}$ of disjoint intervals $[n_j,r(n_j))$ of input lengths such that for each $j$ and $n\in[n_j,r(n_j))$, we have that $s(n)>s_1(n)$.
If there is no $2^{s_2(n)^3}$-robust set $S$ on which $M_2$ solves SAT correctly, there must be a constant $k$ such that for each $n$, there is $m<2^{ks_2(n)^3}$ such that $M_2$ fails to solve SAT correctly on some input of length $m$. Now we use the existence of the sequence $\{I_j\}$ of intervals from the Claim. Since $r(n)=n^{\log(n)}$ is significantly larger than $2^{ks_2(n)^3}$ for large enough $n$, it follows that there is an infinite sequence of input lengths $\{m_i\}$ such that the following hold:(i) $M_2$ fails to solve SAT correctly on some input $x$ of length $m_i$ such that $M$ on input $x$ asks a query with parameter at least $s_1(m_i)$ (ii) For all input lengths $m$ such that $m_i\leq m\leq2^{s_2(m_i)^4}$, $s(m)\geq s_1(m)$.
Now we apply the main idea of~\cite{ref12}, by attempting to use the simulation $M_2$ itself to find a small set of inputs for which $M_2$ fails to decide SAT correctly on at least one of these inputs. We define a Turing machine $A$ as follows. On input $1^m$, $A$ formulates the question of whether there is an instance of size $m$ on which $M_2$ fails to solve SAT correctly and asks a query with parameter at least $s_2(m)$ as an instance $y$ of SAT. Since $M_2$ runs in time $\poly(n)2^{s_2(n)^2}$, such an instance $y$ of size at most $2^{s_2(m)^3}$ for large enough $m$ can be computed in time at most $2^{s_2(m)^3}$. $A$ then runs $M_2$ on $y$ to find a satisfying assignment for $y$, as in the proof of the main result of~\cite{ref12}. For all $m=m_i$ for some $i$, either $A$ succeeds, in which case it finds an instance $x$ of length $m$ on which $M_2$ fails, or it finds a set of at most three instances $y_1,y_2,y_3$ such that $|y_j|=|y|$ for each $1\leq j\leq3$, and $M_2$ fails on at least one of these three instances. Note that by item(ii) in the para above, in the first case $x$ must ask some query with parameter at least $s_1(m)$ that is answered wrongly, and in the second case, one of $y_1,y_2,y_3$ must ask some query with parameter at least $s_1(|y|)$ that is answered wrongly. In either case, a wrongly answered query must be the truth table of a hard function. In the first case, we have that the concatenation $Z_i$ of all queries asked by $x$ is the truth table of a function on $O(\log(m))$ input bits that does not have circuits of size $s_1(m)$, and in the second case we have that the concatenation $Z_i$ of all queries asked by $y_1,y_2,y_3$ must be the truth table of a function on $O(\log(|y|))$ input bits that does not have circuits of size $s_1(|y|)$. $A$ outputs $Z_i$.
In either case, $A$ runs in time $\poly(|Z_i|)2^{s_2(|Z_i|)^3}$ and outputs a string $Z_i$ that does not have circuits of size $s_1(|Z_i|)$. Since $s_2(|Z|)=s_1(|Z|)^{o(1)}$, this implies that $\mathrm E\not\subseteq\SIZE(\poly)$.
\end{proof}

\begin{thebibliography}{24}
\bibitem{ref1}Eric Allender. The complexity of complexity. In Computability and Complexity - Essays Dedicated to Rodney G. Downey on the Occasion of His 60th Birthday, pages 79–94, 2017.
\bibitem{ref2}Eric Allender, Harry Buhrman, Michal Koucký, Dieter van Melkebeek, and Detlef Ronneburger. Power from random strings. SIAM J. Comput., 35(6):1467–1493, 2006.
\bibitem{ref3}Eric Allender and Bireswar Das. Zero knowledge and circuit minimization. In Symposium on Mathematical Foundations of Computer Science (MFCS), pages 25–32, 2014.
\bibitem{ref4}Eric Allender, Joshua A. Grochow, and Cristopher Moore. Graph isomorphism and circuit size. CoRR, abs/1511.08189, 2015. arXiv:1511.08189.
\bibitem{ref5}Eric Allender, Lisa Hellerstein, Paul McCabe, Toniann Pitassi, and Michael E. Saks. Minimizing DNF formulas and AC0 circuits given a truth table. Electronic Colloquium on Computational Complexity (ECCC), 126, 2005.
\bibitem{ref6}Eric Allender and Shuichi Hirahara. New insights on the (non-)hardness of circuit minimization and related problems. In International Symposium on Mathematical Foundations of Computer Science (MFCS), pages 54:1–54:14, 2017.
\bibitem{ref7}Eric Allender, Dhiraj Holden, and Valentine Kabanets. The minimum oracle circuit size problem. In International Symposium on Theoretical Aspects of Computer Science (STACS), pages 21–33, 2015.
\bibitem{ref8}Eric Allender, Michal Koucký, Detlef Ronneburger, and Sambuddha Roy. The pervasive reach of resource-bounded Kolmogorov complexity in computational complexity theory. J. Comput. Syst. Sci., 77(1):14–40, 2011.
\bibitem{ref9}Sanjeev Arora and Boaz Barak. Computational Complexity: A Modern Approach. Cambridge University Press, 1st edition, 2009.
\bibitem{ref10}Sebastian Czort. The complexity of minimizing disjunctive normal form formulas. Master’s Thesis, University of Aarhus, 1999.
\bibitem{ref11}Lance Fortnow and Rahul Santhanam. Robust simulations and significant separations. Infor- mation and Computation, 256:149–159, 2017.
\bibitem{ref12}Dan Gutfreund, Ronen Shaltiel, and Amnon Ta-Shma. If NP languages are hard on the worst- case, then it is easy to find their hard instances. Computational Complexity, 16(4):412–441, 2007.
\bibitem{ref13}Juris Hartmanis and Richard Stearns. On the computational complexity of algorithms. Transactions of the American Mathematical Society, pages 285–306, 1965.
\bibitem{ref14}Shuichi Hirahara, Igor Carboni Oliveira, and Rahul Santhanam. Np-hardness of minimum circuit size problem for OR-AND-MOD circuits. In 33rd Computational Complexity Conference, CCC 2018, pages 5:1–5:31, 2018.
\bibitem{ref15}Shuichi Hirahara and Rahul Santhanam. On the average-case complexity of MCSP and its variants. In Computational Complexity Conference (CCC), pages 7:1–7:20, 2017.
\bibitem{ref16}Shuichi Hirahara and Osamu Watanabe. Limits of minimum circuit size problem as oracle. In Conference on Computational Complexity (CCC), pages 18:1–18:20, 2016.
\bibitem{ref17}John M. Hitchcock and Aduri Pavan. On the NP-completeness of the minimum circuit size problem. In Conference on Foundation of Software Technology and Theoretical Computer Science (FSTTCS), pages 236–245, 2015.
\bibitem{ref18}Valentine Kabanets and Jin-yi Cai. Circuit minimization problem. In Symposium on Theory of Computing (STOC), pages 73–79, 2000.
\bibitem{ref19}Richard Karp and Richard Lipton. Some connections between nonuniform and uniform complexity classes. In Proceedings of the 12th Annual ACM Symposium on Theory of Computing, April 28-30, 1980, Los Angeles, California, USA, pages 302–309, 1980.
\bibitem{ref20}William J. Masek. Some NP-complete set covering problems. Unpublished Manuscript, 1979.
\bibitem{ref21}Cody Murray and Ryan Williams. On the (non) NP-hardness of computing circuit complexity. In Conference on Computational Complexity (CCC), pages 365–380, 2015.
\bibitem{ref22}Igor Carboni Oliveira and Rahul Santhanam. Conspiracies between learning algorithms, circuit lower bounds, and pseudorandomness. In Computational Complexity Conference (CCC), pages 18:1–18:49, 2017.
\bibitem{ref23}Alexander A. Razborov and Steven Rudich. Natural proofs. J. Comput. Syst. Sci., 55(1):24–35, 1997.
\bibitem{ref24}Boris A. Trakhtenbrot. A survey of Russian approaches to perebor (brute-force searches) algorithms. IEEE Annals of the History of Computing, 6(4):384–400, 1984.
\end{thebibliography}
\end{document}
